\subsection{Minkowski 不等式}

设 X 是一个局部紧空间，\(\mu\) 是 X 上的一个测度。在定义于 X 上的正数值函数（有限或非有限）所成的集中，函数 \(|\mu|^{*}(f)\) 是正的、正齐次的、递增的且凸的（ \(\S1\) , No.~3, 命题 10、11 与 12）。

命题 1. --- 对每个实数 \(p \geqslant 1\) 与每对正函数 \(f, g\) （有限或非有限，定义于 \(X\) 上），

\[
\left(| \mu | ^ {*} ((f + g) ^ {p})\right) ^ {1 / p} \leqslant \left(| \mu | ^ {*} (f ^ {p})\right) ^ {1 / p} + \left(| \mu | ^ {*} (g ^ {p})\right) ^ {1 / p}\tag{1}
\]

（Minkowski 不等式）。

当右端的项之一等于 \(+\infty\) 时，不等式 (1) 是明显的。在相反的情形，\(f\) 与 \(g\) 几乎处处有限（§2, No.~3, 命题 7）。若 \(f_1\) 与 \(g_1\) 是分别等价于 \(f\) 与 \(g\) 的有限正函数，则 \(f_1^p\) 、 \(g_1^p\) 与 \((f_1 + g_1)^p\) 分别等价于 \(f^p\) 、 \(g^p\) 与 \((f + g)^p\) ，而由于等价的正函数具有相同的上积分（§2, No.~3, 命题 6），我们只需就 \(f\) 与 \(g\) 处处有限的情形证明不等式 (1) ；但在这一情形，该不等式是 Ch. I, No.~2 的命题 3 中证明的一般 Minkowski 不等式的特例。

我们还将有机会用到下述初等不等式：若 \(p \geqslant 1\) ，则对任意数 \(a \geqslant 0\) 、 \(b \geqslant 0\) ，

\[
a ^ {p} + b ^ {p} \leqslant (a + b) ^ {p}.\tag{2}
\]

若 \(a = b = 0\) ，或数 \(a, b\) 之一等于 \(+\infty\) ，则不等式是明显的；若 \(a, b\) 有限且 \(a + b > 0\) ，则它可写成

\[
\left(\frac {a}{a + b}\right) ^ {p} + \left(\frac {b}{a + b}\right) ^ {p} \leqslant 1,
\]

这由以下事实得出：

\[
\left(\frac {a}{a + b}\right) ^ {p} \leqslant \frac {a}{a + b}, \quad \left(\frac {b}{a + b}\right) ^ {p} \leqslant \frac {b}{a + b} \quad \text { and } \quad \frac {a}{a + b} + \frac {b}{a + b} = 1.
\]

